\documentclass[10pt,a4paper]{amsart}
\usepackage{amsmath,amssymb,amsthm,graphicx}
% For commutative diagrams we use Xy-pic
\usepackage[all]{xy}

% We use 2cell for 2-commutative diagrams.
\xyoption{2cell}
% Theorem environments.
%
\theoremstyle{plain}
\newtheorem{theorem}[subsection]{Theorem}
\newtheorem{proposition}[subsection]{Proposition}
\newtheorem{lemma}[subsection]{Lemma}

\theoremstyle{definition}
\newtheorem{definition}[subsection]{Definition}
\newtheorem{example}[subsection]{Example}
\newtheorem{exercise}[subsection]{Exercise}
\newtheorem{situation}[subsection]{Situation}

\theoremstyle{remark}
\newtheorem{remark}[subsection]{Remark}
\newtheorem{remarks}[subsection]{Remarks}

\numberwithin{equation}{subsection}

% Macros
%
\def\lim{\mathop{\mathrm{lim}}\nolimits}
\def\colim{\mathop{\mathrm{colim}}\nolimits}
\def\Spec{\mathop{\mathrm{Spec}}}
\def\Hom{\mathop{\mathrm{Hom}}\nolimits}
\def\Ext{\mathop{\mathrm{Ext}}\nolimits}
\def\SheafHom{\mathop{\mathcal{H}\!\mathit{om}}\nolimits}
\def\SheafExt{\mathop{\mathcal{E}\!\mathit{xt}}\nolimits}
\def\Sch{\mathit{Sch}}
\def\Mor{\mathop{\mathrm{Mor}}\nolimits}
\def\Ob{\mathop{\mathrm{Ob}}\nolimits}
\def\Sh{\mathop{\mathit{Sh}}\nolimits}
\def\NL{\mathop{N\!L}\nolimits}
\def\CH{\mathop{\mathrm{CH}}\nolimits}
\def\proetale{{pro\text{-}\acute{e}tale}}
\def\etale{{\acute{e}tale}}
\def\QCoh{\mathit{QCoh}}
\def\Ker{\mathop{\mathrm{Ker}}}
\def\Im{\mathop{\mathrm{Im}}}
\def\Coker{\mathop{\mathrm{Coker}}}
\def\Coim{\mathop{\mathrm{Coim}}}

% Boxtimes
%
\DeclareMathSymbol{\boxtimes}{\mathbin}{AMSa}{"02}

%
% Macros for moduli stacks/spaces
%
\def\QCohstack{\mathcal{QC}\!\mathit{oh}}
\def\Cohstack{\mathcal{C}\!\mathit{oh}}
\def\Spacesstack{\mathcal{S}\!\mathit{paces}}
\def\Quotfunctor{\mathrm{Quot}}
\def\Hilbfunctor{\mathrm{Hilb}}
\def\Curvesstack{\mathcal{C}\!\mathit{urves}}
\def\Polarizedstack{\mathcal{P}\!\mathit{olarized}}
\def\Complexesstack{\mathcal{C}\!\mathit{omplexes}}
% \Pic is the operator that assigns to X its picard group, usage \Pic(X)
% \Picardstack_{X/B} denotes the Picard stack of X over B
% \Picardfunctor_{X/B} denotes the Picard functor of X over B
\def\Pic{\mathop{\mathrm{Pic}}\nolimits}
\def\Picardstack{\mathcal{P}\!\mathit{ic}}
\def\Picardfunctor{\mathrm{Pic}}
\def\Deformationcategory{\mathcal{D}\!\mathit{ef}}
\newenvironment{reference}{}{}
\newenvironment{slogan}{}{}
\newenvironment{history}{}{}
\usepackage{cleveref}
\crefname{originallemma}{original lemma}{original lemmas}
\Crefname{originallemma}{Original lemma}{Original lemmas}
\makeatletter
\newlabel{lemma-bound-neighbours}{{0C9U}{source}{}{sourcepointer-1}{}}
\newlabel{lemma-bound-neighbours@cref}{{[originallemma][0][]0C9U}{[1][0][]source}{}{}{}}
\makeatother
\crefname{originallemma}{original lemma}{original lemmas}
\Crefname{originallemma}{Original lemma}{Original lemmas}
\makeatletter
\newlabel{lemma-bound-heart}{{0C9V}{source}{}{sourcepointer-2}{}}
\newlabel{lemma-bound-heart@cref}{{[originallemma][0][]0C9V}{[1][0][]source}{}{}{}}
\makeatother
\crefname{originalproposition}{original proposition}{original propositions}
\Crefname{originalproposition}{Original proposition}{Original propositions}
\makeatletter
\newlabel{proposition-classify-subgraphs}{{0C8Q}{source}{}{sourcepointer-3}{}}
\newlabel{proposition-classify-subgraphs@cref}{{[originalproposition][0][]0C8Q}{[1][0][]source}{}{}{}}
\makeatother
\crefname{originallemma}{original lemma}{original lemmas}
\Crefname{originallemma}{Original lemma}{Original lemmas}
\makeatletter
\newlabel{lemma-Dn}{{0C8D}{source}{}{sourcepointer-4}{}}
\newlabel{lemma-Dn@cref}{{[originallemma][0][]0C8D}{[1][0][]source}{}{}{}}
\makeatother
\crefname{originallemma}{original lemma}{original lemmas}
\Crefname{originallemma}{Original lemma}{Original lemmas}
\makeatletter
\newlabel{lemma-long}{{0C89}{source}{}{sourcepointer-5}{}}
\newlabel{lemma-long@cref}{{[originallemma][0][]0C89}{[1][0][]source}{}{}{}}
\makeatother
\title{In a minimal numerical type of genus g at least two, each weighted intersection number is bounded by seven hundred sixty-eight times g}
\author{The Stacks Project Authors}
\date{stacks-project:models.tex:local-snapshot-20261012; GNU FDL 1.2 or later}
\begin{document}
\maketitle
\noindent\textit{Editorial note.} The counted unit is the Stacks Project Lemma lemma-bound-wm at source lines 4171-4175 together with its proof at source lines 4177-4283, a straight adjacent pairing in the chapter file and one of the longest proofs of this chapter. The chapter is the upstream file models.tex; the pool directory carries the longer name stacks-models and is not used as the chapter name anywhere in this card.\par\medskip
\noindent\textit{Editorial note.} This card has no context span: the statement and proof are self-contained apart from the five pointers, and the notions they use are recorded in the standard prerequisites.\par\medskip
\noindent\textit{Editorial note.} The retained proof names five results of the same chapter: the elementary bound on neighbouring entries, the bound on the heart, the classification of proper subgraphs, and two lemmas about specific shapes of the intersection graph. All five are declared as non-load-bearing navigation pointers with their stable Stacks tags looked up in the shared tag index. Three of them belong to the block of lemmas whose proofs this lane does not count, because those proofs are one-line citations; pointing at their statements from another card is nevertheless valid and none of their intervals is reproduced here.\par\medskip
\noindent\textit{Editorial note.} No citation occurs in the retained spans, so no bibliography entry is supplied.\par\medskip
\noindent\textit{Allowed standard prerequisites (editorial).}\begin{itemize}
\item Minimal numerical types of genus at least two and the classification of their proper subgraphs. A numerical type is a tuple of an integer n and integer families of multiplicities, intersection numbers, weights and genera subject to the five defining conditions of the chapter, and it is minimal when no index has genus zero and self-intersection equal to the negative of its weight; the genus of the type is the alternating sum determined by those data. The chapter classifies the proper subgraphs of the intersection graph of a minimal type of genus at least two, up to the possible shapes, and proves the auxiliary estimates on the entries of that graph that the proof uses. The author's own lemmas and proposition of the same chapter are the allowed inputs.. The counted lemma is the quantitative estimate that makes the classification of minimal types of fixed genus finite, and the chapter's enumeration of minimal types of low genus depends on it.
\end{itemize}\medskip
\ifdefined\hypertarget\hypertarget{sourcepointer-1}{}\fi
\noindent\textit{Source reference.} Original lemma 0C9U: focused-seeds/AGP/sources/stacks-snapshot-20261011/tags.txt, source lines 15363--15363. Navigation pointer invoked by the retained proof at its first line to reduce the assertion to a single entry. This lane does not count that lemma, because its own proof defers the argument to the reader; pointing at its statement from another card is nevertheless valid and its interval is not reproduced. The number is the stable Stacks tag of the target, looked up in the shared tag index.\par\medskip
\ifdefined\hypertarget\hypertarget{sourcepointer-2}{}\fi
\noindent\textit{Source reference.} Original lemma 0C9V: focused-seeds/AGP/sources/stacks-snapshot-20261011/tags.txt, source lines 15364--15364. Navigation pointer invoked by the retained proof for the auxiliary estimate on the entries of the intersection graph; the cited statement is an allowed input of the same section and its proof is not reproduced here. The number is the stable Stacks tag of the target, looked up in the shared tag index.\par\medskip
\ifdefined\hypertarget\hypertarget{sourcepointer-3}{}\fi
\noindent\textit{Source reference.} Original proposition 0C8Q: focused-seeds/AGP/sources/stacks-snapshot-20261011/tags.txt, source lines 15333--15333. Navigation pointer invoked by the retained proof for the classification of the proper subgraphs of the intersection graph; the cited statement is an allowed input of the same section and its proof is not reproduced here. The number is the stable Stacks tag of the target, looked up in the shared tag index.\par\medskip
\ifdefined\hypertarget\hypertarget{sourcepointer-4}{}\fi
\noindent\textit{Source reference.} Original lemma 0C8D: focused-seeds/AGP/sources/stacks-snapshot-20261011/tags.txt, source lines 15302--15302. Navigation pointer invoked by the retained proof in its case analysis for one of the shapes of the intersection graph. That lemma belongs to the block whose proofs this lane does not count, because they are bare citations; pointing at its statement is nevertheless valid and its interval is not reproduced. The number is the stable Stacks tag of the target, looked up in the shared tag index.\par\medskip
\ifdefined\hypertarget\hypertarget{sourcepointer-5}{}\fi
\noindent\textit{Source reference.} Original lemma 0C89: focused-seeds/AGP/sources/stacks-snapshot-20261011/tags.txt, source lines 15298--15298. Navigation pointer invoked by the retained proof in its case analysis for another shape of the intersection graph. Like the previous entry it belongs to the block this lane does not count, and it is pointed at without its interval being reproduced. The number is the stable Stacks tag of the target, looked up in the shared tag index.\par\medskip
\begin{lemma}
\label{lemma-bound-wm}
Fix $g \geq 2$. For every minimal numerical type $n, m_i, a_{ij}, w_i, g_i$
of genus $g$ we have $m_i|a_{ij}| \leq 768g$.
\end{lemma}

\begin{proof}
By Lemma \ref{lemma-bound-neighbours} it suffices to show
$m_i|a_{ii}| \leq 768g$ for all $i$.
Let $J \subset \{1, \ldots, n\}$ be the set of non-$(-2)$-indices
as in Lemma \ref{lemma-bound-heart}. Observe that $J$ is
nonempty as $g \geq 2$. Also $m_j|a_{jj}| \leq 6g$ for $j \in J$
by the lemma.

\medskip\noindent
Suppose we have $j \in J$ and a sequence $i_1, \ldots, i_7$
of $(-2)$-indices such that $a_{ji_1}$ and $a_{i_1i_2}$,
$a_{i_2i_3}$, $a_{i_3i_4}$, $a_{i_4i_5}$, $a_{i_5i_6}$, and
$a_{i_6i_7}$ are nonzero. Then we see from
Lemma \ref{lemma-bound-neighbours}
that $m_{i_1}w_{i_1} \leq 6g$ and $m_{i_1}a_{ji_1} \leq 6g$.
Because $i_1$ is a $(-2)$-index, we have $a_{i_1i_1} = -2w_{i_1}$
and we conclude that $m_{i_1}|a_{i_1i_1}| \leq 12g$.
Repeating the argument we conclude that
$m_{i_2}w_{i_2} \leq 12g$ and $m_{i_2}a_{i_1i_2} \leq 12g$.
Then $m_{i_2}|a_{i_2i_2}| \leq 24g$ and so on.
Eventually we conclude that
$m_{i_k}|a_{i_ki_k}| \leq 2^k(6g) \leq 768g$ for $k = 1, \ldots, 7$.

\medskip\noindent
Let $I \subset \{1, \ldots, n\} \setminus J$ be a maximal connected subset.
In other words, there does not exist a nonempty proper subset
$I' \subset I$ such that
$a_{i'i} = 0$ for $i' \in I'$ and $i \in I \setminus I'$
and $I$ is maximal with this property. In particular, since
a numerical type is connected by definition, we see that there
exists a $j \in J$ and $i \in I$ with $a_{ij} > 0$.
Looking at the classification
of such $I$ in Proposition \ref{proposition-classify-subgraphs}
and using the result of the previous paragraph, we see that
$w_i|a_{ii}| \leq 768g$ for all $i \in I$ unless $I$ is as described in
Lemma \ref{lemma-long} or
Lemma \ref{lemma-Dn}.
Thus we may assume the nonvanishing of $a_{ii'}$, $i, i' \in I$
has either the shape
$$
\xymatrix{
\bullet \ar@{-}[r] &
\bullet \ar@{-}[r] &
\bullet \ar@{..}[r] &
\bullet \ar@{-}[r] &
\bullet \ar@{-}[r] &
\bullet
}
$$
(which has 3 subcases as detailed in Lemma \ref{lemma-long})
or the shape
$$
\xymatrix{
\bullet \ar@{-}[r] &
\bullet \ar@{-}[r] &
\bullet \ar@{..}[r] &
\bullet \ar@{-}[r] &
\bullet \ar@{-}[r] \ar@{-}[d] &
\bullet \\
& & & & \bullet
}
$$
We will prove the bound holds for the first subcase of
Lemma \ref{lemma-long} and leave the other cases to reader (the argument
is almost exactly the same in those cases).

\medskip\noindent
After renumbering we may assume $I = \{1, \ldots, t\} \subset \{1, \ldots, n\}$
and there is an integer $w$ such that
$$
w = w_1 = \ldots = w_t =
a_{12} = \ldots = a_{(t - 1) t} =
-\frac{1}{2} a_{i_1i_2} = \ldots = -\frac{1}{2} a_{(t - 1) t}
$$
The equalities $a_{ii}m_i + \sum_{j \not = i} a_{ij}m_j = 0$ imply
that we have
$$
2m_2 \geq m_1 + m_3, \ldots, 2m_{t - 1} \geq m_{t - 2} + m_t
$$
Equality holds in $2m_i \geq m_{i - 1} + m_{i + 1}$
if and only if $i$ does not ``meet'' any indices besides
$i - 1$ and $i + 1$. And if $i$ does meet
another index, then this index is in $J$ (by maximality of $I$).
In particular, the map
$\{1, \ldots, t\} \to \mathbf{Z}$, $i \mapsto m_i$ is concave.

\medskip\noindent
Let $m = \max(m_i, i \in \{1, \ldots, t\})$. Then
$m_i|a_{ii}| \leq 2mw$ for $i \leq t$ and
our goal is to show that $2mw \leq 768g$.
Let $s$, resp.\ $s'$ in $\{1, \ldots, t\}$ be the
smallest, resp.\ biggest index with $m_s = m = m_{s'}$.
By concavity we see that $m_i = m$ for $s \leq i \leq s'$.
If $s > 1$, then we do not have equality in
$2m_s \geq m_{s - 1} + m_{s + 1}$
and we see that $s$ meets an index from $J$.
In this case $2mw \leq 12g$ by the result of the second paragraph
of the proof.
Similarly, if $s' < t$, then $s'$ meets an index from $J$
and we get $2mw \leq 12g$ as well.
But if $s = 1$ and $s' = t$, then we conclude
that $a_{ij} = 0$ for all $j \in J$ and $i \in \{2, \ldots, t - 1\}$.
But as we've seen that there must be a pair $(i, j) \in I \times J$
with $a_{ij} > 0$, we conclude that this happens either with
$i = 1$ or with $i = t$ and we conclude $2mw \leq 12g$
in the same manner as before (as $m_1 = m = m_t$ in this case).
\end{proof}

\end{document}
